%%%%%%%%%%%%%%%%%%%%%%%%%%%%%%%%%%%%%%%%%%%%%%%%%%%%%%%%%%%%%%
%%%%%%%%%%%% MA-0701 Seminario de Matemática %%%%%%%%%%%%%%%%%%
%%%%%%%%%%%%%%%%%%%%%%%%%%%%%%%%%%%%%%%%%%%%%%%%%%%%%%%%%%%%%%
\newpage
\subsection*{MA-0701 Seminario de Matemática}
\addcontentsline{toc}{subsection}{MA-0701 Seminario de Matemática}

\hypertarget{MA0701}{}
\begin{refsection}
\begin{table}[H]
\centering{
\begin{tabular}{|ll|}
\hline
\textbf{Año:} IV  & \textbf{Requisitos:} Haber aprobado al menos uno de: MA-0661 \\ & Algebra Abstracta II, MA-0635 Introducción a la Topología \\ & o MA-0641 Matemática Computacional II \\ 
\textbf{Ciclo:} VII  & \textbf{Correquisitos:} Ninguno \\ 
\textbf{Tipo de curso:} Seminario & \textbf{Horario:} Según calendario del Coloquio \\ & del Departamento de Matemáticas \\ 
\textbf{Créditos:} 0 & \textbf{Modalidad:} Presencial \\ \hline
\end{tabular}
}
\end{table}

\subsubsection*{Descripción del curso}

Este curso está dirigido a personas estudiantes de los últimos ciclos de la carrera Bachillerato en Matemática. Ofrece un espacio de formación complementaria mediante la asistencia a charlas, coloquios, seminarios y conferencias sobre temas relevantes de la matemática contemporánea y de su historia, organizados o avalados por el Departamento de Matemáticas.

A diferencia de un curso teórico tradicional, la actividad central consiste en participar en el ciclo de exposiciones del departamento y en una sesión final de síntesis, en la cual se analizan los contenidos vistos y los aspectos formales de la comunicación matemática. El propósito es acercar al estudiantado a la práctica profesional e investigativa de la disciplina, ampliar su panorama sobre áreas activas de la matemática y fortalecer su capacidad de escuchar, interpretar y discutir exposiciones especializadas.

\subsubsection*{Objetivos}

\textbf{General:} Que la persona estudiante sea capaz de analizar e integrar criterios para valorar de manera fundamentada exposiciones científicas en el ámbito de las matemáticas.

\textbf{Específicos:} Al finalizar el curso se espera que la persona estudiante sea capaz de:

\begin{enumerate}
% \item Asistir de manera responsable a charlas, coloquios o conferencias matemáticas, registrando su participación conforme a las normas del curso.

\item Identificar y describir temas, problemas y líneas de investigación actuales o clásicos abordados en las exposiciones.

\item Relacionar las charlas asistidas con conocimientos previos de la malla curricular, estableciendo conexiones coherentes y significativas.

\item Analizar e integrar, de forma fundamentada, aspectos del contenido y de la forma de las exposiciones (claridad, estructura, calidad de la presentación, fluidez, manejo del público y capacidad de comunicación).

\item Sintetizar y argumentar críticamente ideas en la sesión final del semestre, generando reflexiones transferibles sobre la comunicación científica en matemáticas.
\end{enumerate}

\subsubsection*{Contenidos}

El contenido del curso no sigue un temario fijo de lecciones magistrales, sino el calendario de actividades válidas que el Departamento de Matemáticas anuncie para cada ciclo lectivo. Entre los tipos de temas que pueden abordarse se incluyen, a título de ejemplo:

\begin{itemize}
\item Avances y problemas abiertos en distintas áreas de la matemática pura o aplicada.
\item Historia y desarrollo de ideas, teoremas o escuelas matemáticas.
\item Matemática interdisciplinaria y aplicaciones en otras ciencias o en la industria.
\item Comunicación, divulgación y enseñanza de las matemáticas.
\item Metodología de investigación y presentación de resultados.
\end{itemize}



\subsubsection*{Metodología}

El curso adopta la modalidad \textbf{seminario}: la persona estudiante participa de forma activa asistiendo a las actividades programadas y concurriendo a la sesión final de discusión. La coordinación del curso prepara o valida el plan preliminar de actividades del ciclo lectivo y divulga oportunamente fechas, lugares e instrucciones específicas.

El curso comprende:

\begin{enumerate}
\item \textbf{Asistencia a actividades válidas (transversal al ciclo lectivo):}
\begin{enumerate}
    \item Charlas del Coloquio del Departamento de Matemáticas.
    \item Seminarios, conferencias o ciclos especiales anunciados por el departamento como válidos para este curso.
    \item Otras actividades académicas que la coordinación del curso disponga expresamente, con indicaciones oportunas.
\end{enumerate}

\item \textbf{Sesión final de síntesis y discusión (1 sesión, fin de semestre):}
\begin{enumerate}
    \item Recapitulación guiada de las charlas asistidas durante el año.
    \item Discusión sobre el contenido matemático y sobre la forma de las exposiciones: claridad, calidad de la presentación, fluidez al exponer, capacidad de comunicación, manejo del público y estructuración de la charla.
\end{enumerate}
\end{enumerate}

\noindent\textbf{Lineamientos metodológicos:} La persona docente coordinadora del curso debe:

\begin{enumerate}
    \item Publicar al inicio del ciclo lectivo un cronograma tentativo de actividades válidas y confirmar cada evento con antelación razonable.
    \item Llevar o supervisar un sistema de control de asistencia institucionalmente acordado.
    \item Orientar al estudiantado sobre las expectativas del curso, las normas de cumplimiento y los criterios de evaluación.
    \item Dirigir la sesión final de síntesis, favoreciendo la participación, el respeto de opiniones diversas y la discusión fundamentada.
    \item Coordinar, orientar y evaluar la discusión grupal, sintetizando los argumentos tratados.
    \item Salvo casos que lo ameriten, la persona coordinadora debe ser quien ocupe la dirección del Departamento de Matemáticas, o quien este designe.
\end{enumerate}

\textbf{Sugerencias metodológicas:}

\begin{enumerate}

    \item Pueden validarse charlas de ciclos especiales del departamento (por ejemplo, cátedras, semanas temáticas o invitados internacionales), siempre que se anuncien con claridad como actividades válidas para MA-0701.
    \item Se recomienda solicitar a las personas estudiantes breves notas reflexivas sobre algunas charlas, como insumo para la sesión final; ello queda a criterio de la coordinación.
    \item El horario de matrícula del curso coincide con el horario de las actividades del Coloquio o de las sesiones válidas anunciadas; no implica horas fijas semanales de aula magna.
\end{enumerate}

\noindent\textbf{Normas de cumplimiento:}

Para que una actividad cuente como una asistencia válida, como mínimo debe cumplirse lo siguiente:

\begin{itemize}
\item Que el Departamento de Matemáticas o la coordinación del curso la haya dispuesto expresamente como válida para MA-0701.
\item Firmar la lista de asistencia (o el mecanismo equivalente aprobado por el departamento).
\item Llegar puntualmente; una vez retirada la lista para firma no se contará la asistencia.
\item Permanecer hasta el final de la actividad.
\item Asistencia presencial, salvo que las instrucciones del evento autoricen explícitamente otra modalidad.
\item Cumplir, cuando corresponda, condiciones adicionales comunicadas con anticipación (por ejemplo, registro previo o encuesta posterior).
\end{itemize}

%

\subsubsection*{Evaluación}

El curso se califica en escala \textbf{aprobado / reprobado}.

\begin{enumerate}
\item \textbf{Asistencia mínima:} aprobar el curso requiere acreditar un mínimo de \textbf{10 asistencias} a actividades válidas. El Departamento de Matemáticas garantizará anunciar suficientes actividades en cada ciclo lectivo para que el estudiantado pueda cumplir este requisito.
\item \textbf{Participación en la sesión final:} la asistencia y participación activa en la sesión de cierre del semestre es obligatoria. Las ausencias podrán justificarse conforme al Reglamento del Régimen Académico Estudiantil y demás normativa institucional vigente, a criterio de la coordinación del curso.
\item \textbf{Plazo para completar las asistencias:} la persona estudiante dispone de \textbf{dos ciclos lectivos consecutivos} para completar las 10 asistencias, matriculándose una sola vez al inicio de ese periodo:
\begin{itemize}
\item Quien matricula en el \textbf{I ciclo} de un año dispone del I y II ciclo de ese mismo año para finalizar; no debe volver a matricular el curso en el II ciclo de ese año.
\item Quien matricula en el \textbf{II ciclo} de un año dispone del II ciclo de ese año y del I ciclo del año siguiente; no debe volver a matricular el curso en el I ciclo del año siguiente.
\end{itemize}
\item \textbf{No aprobación:} si no se completan las 10 asistencias en el plazo correspondiente, o si falta a la sesión final sin justificación válida, el curso se reportará como \textbf{No Aprobado (NAP)} y deberá matricularse nuevamente, reiniciando el conteo de asistencias desde cero.
\end{enumerate}

\subsubsection*{Bibliografía}

No aplica de manera fija. Cada charla o conferencia puede indicar referencias propias, las cuales serán comunicadas por el expositor o la coordinación del evento.

\end{refsection}
